\chapter{Cycle Detection, Checkpointing, and Data Staging Modes}
\label{ch:lessons-05-07}

\begin{lessonbox}{Lecture 09 --- Validating graphs and cycles}
\textbf{Objective.} Learn how DAGonStar rejects invalid workflow graphs.

\textbf{Concept.} A DAG cannot contain cycles. If \texttt{A} depends on
\texttt{B} and \texttt{B} depends on \texttt{A}, no valid execution order
exists.

\textbf{Course example} (\texttt{examples/09-cycles/cycles.py}, Lecture~09).
\lstinputlisting[caption={Lecture 09: \texttt{cycles.py}},label={lst:lecture-09}]{../../examples/09-cycles/cycles.py}

\textbf{Run} from the repository root:
\begin{lstlisting}[style=bash]
python examples/09-cycles/cycles.py
\end{lstlisting}

\textbf{Recorded output} (\texttt{examples/09-cycles/output/cycles.txt}):
\lstinputlisting[style=bash]{../../examples/09-cycles/output/cycles.txt}

\textbf{Scientific practice note.} Cycle detection is not only a software
constraint. It also catches ambiguous scientific logic, such as attempting
to calibrate a model using results that the same model has not produced yet.
\end{lessonbox}

\section{Discussion: Cycles as Scientific Modeling Errors}

\Cref{sec:dag-theory} gave the graph-theoretic definition of acyclicity;
Lecture~09 demonstrates the engine enforcing it operationally, via
\texttt{Validate\_WF()}. The scientific practice note deserves emphasis
because it reframes a software constraint as a \emph{modeling} constraint: a
cycle in a workflow DAG is frequently the symptom of a circular scientific
argument --- calibrating a model on its own not-yet-produced output, or
post-processing a product that the post-processing step itself is meant to
generate. Treat a cycle-validation failure as an invitation to re-examine the
experimental design, not merely a bug to silence.

\begin{figure}[h]
\centering
\begin{tikzpicture}[
  node distance=26mm,
  task/.style={rectangle, rounded corners, draw=red!60!black, very thick,
    fill=red!5, minimum width=20mm, minimum height=10mm, font=\sffamily\bfseries},
  arr/.style={-{Stealth[length=2.6mm]}, thick, red!60!black}
]
  \node[task] (A) {A};
  \node[task, right=of A] (B) {B};
  \draw[arr] (A.north east) to[bend left=35] node[above,font=\sffamily\scriptsize]{B depends on A} (B.north west);
  \draw[arr] (B.south west) to[bend left=35] node[below,font=\sffamily\scriptsize]{A depends on B} (A.south east);
\end{tikzpicture}
\caption{An invalid two-node cycle: neither task can legally run first.
\texttt{Validate\_WF()} rejects this graph before execution begins.}
\label{fig:cycle}
\end{figure}

\begin{lessonbox}{Lecture 15 --- Checkpointing}
\textbf{Objective.} Record task completion metadata and resume a workflow.

\textbf{Concept.} \texttt{Workflow(checkpoint\_file=...)} writes checkpoint
metadata after execution. \texttt{workflow.run(resume\_checkpoint\_file=...)}
can reuse completed task directories when the checkpoint records success and
the working directory still exists.

\textbf{Course example} (\texttt{examples/15-checkpoint/checkpoint\_resume.py}, Lecture~15).
\lstinputlisting[caption={Lecture 15: \texttt{checkpoint\_resume.py}},label={lst:lecture-15}]{../../examples/15-checkpoint/checkpoint_resume.py}

\textbf{Run} from the repository root:
\begin{lstlisting}[style=bash]
python examples/15-checkpoint/checkpoint_resume.py
\end{lstlisting}

\textbf{Recorded output} (\texttt{examples/15-checkpoint/output/checkpoint\_resume.txt}):
\lstinputlisting[style=bash]{../../examples/15-checkpoint/output/checkpoint_resume.txt}

\textbf{Scientific practice note.} Checkpointing is most valuable for
expensive simulations. Preserve checkpoint files together with enough
metadata to know which code and configuration produced them.
\end{lessonbox}

\section{Discussion: Checkpoint Semantics and Their Limits}

The checkpoint file written by Lecture~15 has the shape:

\begin{lstlisting}[caption={Representative checkpoint JSON shape.}]
{
  "_scratch_dir": "/tmp/",
  "Experiment.A": {
    "working_dir": "/tmp/1710000000000-A",
    "workflow": "Experiment",
    "name": "A",
    "code": 0
  }
}
\end{lstlisting}

Checkpoint keys are built from \texttt{<workflow-name>.<task-name>}; stable
naming is therefore essential, and renaming a task is, from the checkpoint
system's perspective, indistinguishable from introducing a new task.
Checkpoint reuse additionally requires that the \emph{recorded working
directory still exist} on disk --- a checkpoint entry alone cannot recover
output that scratch-directory garbage collection or manual cleanup has
already removed.

\begin{cautionbox}
Checkpoints record task \emph{success} and a working-directory path, not full
semantic equivalence of code, inputs, or parameters. If a task's command
changes but its name and previous success remain recorded, checkpoint reuse
may incorrectly skip recomputation. For rigorous scientific reproducibility,
store code version, configuration, input checksums, and software environment
\emph{alongside} the checkpoint file, not only inside it.
\end{cautionbox}

A complete reference treatment of checkpoint semantics --- including remote,
Docker, and cloud checkpointing caveats --- is reproduced in Appendix~D.

\begin{lessonbox}{Lecture 13 --- Staging modes}
\textbf{Objective.} Control how DAGonStar stages files between tasks.

\textbf{Concept.} The workflow default data mover is set with
\texttt{workflow.set\_data\_mover(...)}. Local workflows commonly use
\texttt{DataMover.COPY} or \texttt{DataMover.LINK}.

\textbf{Course example} (\texttt{examples/13-staging/staging\_modes.py}, Lecture~13).
\lstinputlisting[caption={Lecture 13: \texttt{staging\_modes.py}},label={lst:lecture-13}]{../../examples/13-staging/staging_modes.py}

\textbf{Run} from the repository root:
\begin{lstlisting}[style=bash]
python examples/13-staging/staging_modes.py
\end{lstlisting}

\textbf{Recorded output} (\texttt{examples/13-staging/output/staging\_modes.txt}):
\lstinputlisting[style=bash]{../../examples/13-staging/output/staging_modes.txt}

\textbf{Scientific practice note.} \texttt{COPY} is often better for archival
reproducibility because downstream tasks get their own materialized inputs.
\texttt{LINK} can reduce storage and I/O cost when intermediate files remain
stable.
\end{lessonbox}

\section{Discussion: Choosing a Staging Policy}

The choice between \texttt{COPY} and \texttt{LINK} is a genuine engineering
trade-off, not a stylistic preference, and it should be made deliberately per
workflow (or even per task, since \texttt{set\_data\_mover} may also be
invoked at task scope):

\begin{itemize}[leftmargin=1.4em]
  \item \textbf{\texttt{COPY}} materializes an independent file in the
    consumer's scratch directory. This maximizes archival reproducibility ---
    a consumer task directory, once created, is self-contained even if the
    producer's directory is later cleaned up --- at the cost of doubled
    storage and I/O for large scientific products (e.g.\ multi-gigabyte
    NetCDF restart files).
  \item \textbf{\texttt{LINK}} creates a symbolic link rather than a copy.
    This is appropriate when intermediate files are large, numerous, and
    known to remain stable for the lifetime of the workflow run, but it
    \emph{breaks} the self-containment property: deleting the producer's
    scratch directory invalidates every consumer that linked to it.
\end{itemize}

\begin{practicebox}
A useful operational rule: use \texttt{COPY} as the default for any
workflow whose outputs will be archived or published, and consider
\texttt{LINK} only for large, purely intermediate products that are deleted
at the end of a single run and never separately archived.
\end{practicebox}

\section{Cycle Detection as an Instance of a General Algorithmic Problem}

\Cref{sec:topological-ordering} already established the precise mathematical
relationship between acyclicity and the existence of a valid execution
order; it is worth, in this chapter, connecting that result to the concrete
\emph{algorithmic} problem \texttt{Validate\_WF()} must solve, since the
choice of algorithm has direct consequences for how the engine scales to
large, realistic scientific workflows.

Cycle detection in a directed graph is classically solved by depth-first
search, tracking, for each node, whether it is unvisited, currently on the
recursion stack (``in progress''), or fully processed (``done''). A cycle is
detected precisely when the search encounters an edge leading to a node
already marked ``in progress'' --- that is, a node reachable from itself
through the edges currently being traversed. This algorithm runs in time
linear in the number of nodes plus edges, $O(|V| + |E|)$, which matters in
practice: a workflow with hundreds or thousands of tasks --- not unrealistic
for an ensemble of environmental simulations, each itself decomposed
according to the granularity discussion of \Cref{sec:task-granularity} ---
can still be validated in time proportional to its own size, rather than,
say, quadratic time that would make validation itself a bottleneck on large
DAGs.

\begin{practicebox}
The practical lesson for a workflow author is reassuring rather than
cautionary: because cycle detection is cheap relative to the cost of running
even a single scientific task, there is essentially no performance reason to
ever omit a \texttt{Validate\_WF()} call, or to avoid restructuring a large
workflow's dependencies during development for fear of expensive
revalidation. The discipline recommended throughout this book --- validate
before running, especially for workflows with expensive or non-idempotent
tasks (\Cref{ch:lessons-01-02}) --- is essentially free in computational
terms.
\end{practicebox}

\section{Checkpointing in the Broader Context of Fault-Tolerant Computing}

Checkpointing is a technique with a long history in fault-tolerant and
high-performance computing, predating workflow engines entirely: long-running
parallel simulations have, for decades, periodically written their internal
state to disk precisely so that a hardware failure partway through a
multi-day run need not force a complete restart from the beginning. It is
useful to situate DAGonStar's task-level checkpointing within this broader
tradition, both to clarify what it does provide and to be precise about
what it does not.

Classical \emph{application-level} checkpointing, of the kind a simulation
code implements internally, captures fine-grained numerical state --- the
full field values of a partially completed timestep, for instance --- and
typically requires the application itself to be written with checkpoint
support in mind. DAGonStar's checkpointing, by contrast, operates entirely
at the \emph{task} granularity: it records whether a task, treated as an
opaque unit, completed successfully, and whether its output directory still
exists. This is a coarser, but considerably more general, form of
checkpointing --- it requires no cooperation from the external program being
wrapped, applies uniformly to a shell command, a Slurm job, a Docker
container, or an LLM completion alike, and composes naturally with the
granularity choices a workflow author makes when decomposing a pipeline
into tasks (\Cref{sec:task-granularity}).

This coarseness has a direct, quantifiable consequence for how a workflow
should be decomposed if checkpoint-driven resumability is a priority for a
given pipeline: the \emph{finer} the task granularity, the more checkpoint
boundaries exist, and the less recomputation a resumed run must perform
after a failure. A single, coarse-grained task that internally runs an
entire multi-hour simulation chain provides exactly one checkpoint boundary
--- success or failure of the whole chain --- whereas the same chain
decomposed into five sequential tasks, each individually checkpointed,
allows a resumed run to skip however many of the five had already
succeeded before the failure occurred. This is a second, independent
argument --- alongside the traceability argument already given in
\Cref{sec:task-granularity} --- for preferring finer-grained task
decomposition in workflows expected to run long enough that mid-run failure
is a realistic possibility.

\section{Staging Cost as a Function of Data Volume}

\Cref{ch:lessons-05-07}'s discussion of \texttt{COPY} versus \texttt{LINK}
staging modes can be made more quantitative by considering explicitly how
each mode's cost scales with the size of the data product being staged.

\texttt{COPY}'s cost --- in both wall-clock time and storage --- scales
linearly with file size: copying a ten-gigabyte NetCDF restart file takes
meaningfully longer, and consumes meaningfully more disk space, than copying
a ten-kilobyte CSV summary. \texttt{LINK}'s cost, by contrast, is
essentially constant regardless of the referenced file's size, since a
symbolic link is a small filesystem object containing only a path string,
independent of the size of the file it points to. For the small,
text-oriented examples used throughout Lectures~06--15 of the course, this
distinction is invisible --- a kilobyte-scale CSV copies essentially
instantly either way --- but it becomes the dominant practical
consideration the moment a workflow scales to the gigabyte-scale data
products characteristic of the environmental modeling case study of
\Cref{ch:case-study}: a multi-stage WRF pipeline that copies, rather than
links, every intermediate NetCDF file between \texttt{geogrid},
\texttt{ungrib}, \texttt{metgrid}, and \texttt{real} could plausibly spend
a non-trivial fraction of total wall-clock time on staging alone, an
overhead \texttt{LINK} would largely eliminate at the cost of the
archival self-containment \Cref{ch:lessons-05-07} already discussed.

\begin{exercise}
For the three-stage \texttt{download} $\to$ \texttt{preprocess} $\to$
\texttt{simulate} checkpoint example discussed earlier in this chapter,
sketch (in prose, no code required) how the choice between \texttt{COPY} and
\texttt{LINK} for the \texttt{download}~$\to$~\texttt{preprocess} edge would
interact with checkpoint reuse after a \texttt{simulate} failure, under the
assumption that \texttt{download}'s output is a forty-gigabyte
satellite-observation file. Which staging mode would you recommend for this
specific edge, and why?
\end{exercise}

\section{Chapter Summary}

Part~II opens with the three mechanisms that make DAGonStar workflows robust
in practice, not merely correct on paper: \texttt{Validate\_WF()} rejects
structurally invalid (cyclic) graphs before execution, using a classical,
linear-time graph algorithm that imposes essentially no practical
performance cost; checkpoint files allow expensive workflows to resume past
completed work, with explicitly documented limits on what ``completed''
guarantees, situated here within the broader, decades-old tradition of
fault-tolerant checkpointing in high-performance computing; and the
\texttt{DataMover} abstraction lets a workflow author trade storage/I-O cost
against archival self-containment, per dependency edge if needed, a
trade-off whose practical weight scales directly with the size of the
scientific data products involved. The rest of the book extends this same
dependency, staging, and lifecycle infrastructure across containers, remote
HPC schedulers, meta-workflows, and --- centrally for this book's stated
subject --- AI and web services.
